\documentclass[aspectratio=169,spanish]{beamer}

% =========================================================
% PAQUETES
% =========================================================

\usepackage[utf8]{inputenc}
\usepackage[spanish, mexico]{babel}
\usepackage{graphicx}
\usepackage{tikz}
\usepackage{booktabs}
\usepackage{array}
\usepackage{ragged2e}
\usepackage{amsmath}
\usepackage{xcolor}
\usepackage{csquotes}

% Bibliografía
\usepackage[
    style=iso-authoryear,
    currentlang
]{biblatex}

\addbibresource{../bibliography/referencias.bib}

\usetikzlibrary{
    arrows.meta,
    positioning,
    calc,
    shapes.geometric,
    fit,
    backgrounds
}

% =========================================================
% RUTAS
% =========================================================

\newcommand{\marcoimg}{../images/marcoref}
\newcommand{\imgdir}{../images}

% =========================================================
% TEMA
% =========================================================

\usetheme{Pittsburgh}

% =========================================================
% COLORES
% =========================================================

\definecolor{UCBBlue}{RGB}{28,73,122}
\definecolor{LightGray}{RGB}{245,246,248}
\definecolor{DarkGray}{RGB}{55,55,55}
\definecolor{SoftBlue}{RGB}{224,233,243}

\setbeamercolor{normal text}{
    fg=DarkGray,
    bg=white
}

\setbeamercolor{structure}{
    fg=UCBBlue
}

\setbeamercolor{frametitle}{
    fg=UCBBlue,
    bg=white
}

\setbeamerfont{title}{
    size=\Large,
    series=\bfseries
}

\setbeamerfont{frametitle}{
    size=\Large,
    series=\bfseries
}

% =========================================================
% NAVEGACIÓN
% =========================================================

\setbeamertemplate{navigation symbols}{}
\setbeamertemplate{footline}{}

% =========================================================
% ENCABEZADO
% Sección izquierda / página derecha
% =========================================================

\setbeamertemplate{headline}{%
    \leavevmode
    \vspace{0.12cm}

    \hbox to \paperwidth{%
        \hspace{0.40cm}%
        {\small\bfseries\color{UCBBlue}\insertsectionhead}%
        \hfill
        {\small\color{DarkGray}
        \insertframenumber/\inserttotalframenumber}%
        \hspace{0.40cm}%
    }

    \vspace{0.10cm}

    {\color{UCBBlue!25}
    \rule{\paperwidth}{0.4pt}}

    \vspace{0.05cm}
}

% =========================================================
% PORTADA
% =========================================================

% A-001: el título definitivo continúa NO DECIDIDO.
% Se conserva aquí la formulación actualmente utilizada
% en perfil.tex hasta decisión académica posterior.

\title{
DISEÑO DE UNA DOSIFICADORA GRAVIMÉTRICA AUTOMATIZADA\\
DE BAJO COSTO PARA HARINA EN PANADERÍAS ARTESANALES
}

\subtitle{
Caso de estudio: Panadería San Miguel
}

\author{
Alesandro Paredes Zenteno
}

\institute{
Universidad Católica Boliviana ``San Pablo''\\
Ingeniería Mecatrónica
}

\date{
Septiembre de 2026
}

\titlegraphic{
    
\includegraphics[
        height=0.17\textheight
    ]{\imgdir/ucbLogo1.png}
}

% =========================================================
% DOCUMENTO
% =========================================================

\begin{document}

% =========================================================
% PORTADA
% =========================================================

% BLOQUES 1 Y 2: apertura, contexto y diagnóstico (diapositivas 1--8).
% Guion de exposición en comentarios; no constituye una decisión académica.
\begin{frame}[plain]
    \maketitle
    % Guion: presentar el caso. El título y el criterio de bajo costo siguen abiertos (A-001).
\end{frame}

\section{Introducción}

% 2 — RECORRIDO
\begin{frame}{Recorrido del Marco Referencial}
\vspace{0.25cm}
\begin{columns}[T,onlytextwidth]
\begin{column}{0.46\textwidth}
{\color{UCBBlue}\Large\bfseries 01\quad Contexto}
\par\medskip
Panadería San Miguel y ubicación de la dosificación de harina.
\par\vspace{0.55cm}
{\color{UCBBlue}\Large\bfseries 02\quad Procedimiento actual}
\par\medskip
Secuencia manual, traslados, pesaje y levantamiento de campo.
\end{column}
\begin{column}{0.46\textwidth}
{\color{UCBBlue}\Large\bfseries 03\quad Evidencia y problema}
\par\medskip
Resultados registrados, proyecciones y planteamiento del problema.
\par\vspace{0.55cm}
{\color{UCBBlue}\Large\bfseries 04\quad Objetivos y delimitación}
\par\medskip
Objetivos, justificación y fronteras del proyecto.
\end{column}
\end{columns}
\vfill
{\small El diagnóstico sustenta el planteamiento del problema.}
% Guion: explicar el recorrido en una frase y pasar al proceso general.
\end{frame}

% 3 — PROCESO GENERAL
\begin{frame}{La dosificación dentro del proceso de panificación}
\centering
\vspace{0.2cm}
\resizebox{0.97\textwidth}{!}{%
    \begin{tikzpicture}[
    node distance=8mm and 8mm,
    >=Latex,
    block/.style={
        rectangle,
        draw,
        minimum height=1.2cm,
        text width=3.2cm,
        align=center,
        font=\small
    },
    alcance/.style={
        block,
        dashed,
        very thick,
        font=\small\bfseries
    },
    arrow/.style={
        -{Latex[width=8pt,length=6pt]},
        thick
    }
]

% -----------------------------------------------------------------
% Primera fila
% -----------------------------------------------------------------

\node (preparacion) [block]
{Preparación de\\ingredientes};

\node (dosificacion) [alcance, right=of preparacion]
{Dosificación de\\ingredientes\\
\textit{\footnotesize Alcance: harina}};

\node (mezclado) [block, right=of dosificacion]
{Mezclado y\\amasado};

\node (formado) [block, right=of mezclado]
{División y\\formado};


% -----------------------------------------------------------------
% Segunda fila
% -----------------------------------------------------------------

\node (fermentacion) [block, below=12mm of formado]
{Fermentación};

\node (horneado) [block, left=of fermentacion]
{Horneado};

\node (producto) [block, left=of horneado]
{Producto\\terminado};


% -----------------------------------------------------------------
% Conexiones
% -----------------------------------------------------------------

\draw [arrow] (preparacion) -- (dosificacion);
\draw [arrow] (dosificacion) -- (mezclado);
\draw [arrow] (mezclado) -- (formado);

\draw [arrow]
    (formado.south)
    -- ++(0,-5mm)
    -| (fermentacion.north);

\draw [arrow] (fermentacion) -- (horneado);
\draw [arrow] (horneado) -- (producto);

\end{tikzpicture}}
\par\vspace{0.4cm}
{\large\color{UCBBlue}\bfseries Dosificación de harina previa al mezclado y amasado}
\par\smallskip
{\small El recuadro discontinuo identifica la etapa estudiada.}
\vfill
\raggedright
{\footnotesize Fuente: Adaptado de \textcite{lim2019spc}.}
% Guion: recorrer el flujo y detenerse en harina. Las operaciones posteriores
% contextualizan el proceso y quedan fuera del desarrollo principal (D-002).
% Verificación pendiente: contraste de la adaptación con el artículo original.
\end{frame}

% 4 — CASO DE ESTUDIO
\begin{frame}{Dosificación de harina en la Panadería San Miguel}
{\small Santa Cruz de la Sierra. Caso de estudio del proyecto.}
\par\vspace{0.2cm}
\begin{columns}[T,onlytextwidth]
\begin{column}{0.48\textwidth}
\centering
\includegraphics[width=\linewidth,height=0.46\textheight,keepaspectratio]
    {\marcoimg/ubicacion_bolsas_harina.jpg}
\par\smallskip
{\small\bfseries Almacenamiento de harina}
\end{column}
\begin{column}{0.48\textwidth}
\centering
\includegraphics[width=\linewidth,height=0.46\textheight,keepaspectratio]
    {\marcoimg/mesa_ingredientes_balanza.jpg}
\par\smallskip
{\small\bfseries Punto de pesaje}
\end{column}
\end{columns}
\vfill
{\large El operario transporta, incorpora y pesa la harina manualmente.}
\par\medskip
{\footnotesize Fuente: Registro fotográfico propio.}
% Guion: identificar ambas zonas. La condición manual describe el procedimiento;
% el diagnóstico cuantitativo se presenta en las diapositivas de resultados.
\end{frame}

\section{Planteamiento del problema}

% 5 — FLUJO GENERAL DEL PROCEDIMIENTO

\begin{frame}{Procedimiento actual de dosificación}

\centering

\begin{tikzpicture}[
  etapa/.style={
    draw=UCBBlue,
    rounded corners=2pt,
    text width=3.75cm,
    minimum height=1.6cm,
    align=center,
    font=\footnotesize
  },
  arr/.style={
    -{Latex[length=2mm]},
    thick,
    UCBBlue
  }
]

% Primera fila
\node[etapa] (a)
{\textbf{1. Preparación y tara}\\
Colocar el recipiente receptor\\
y establecer el cero};

\node[etapa, right=0.75cm of a] (b)
{\textbf{2. Traslado y carga}\\
Recoger, transportar\\
y descargar harina};

\node[etapa, right=0.75cm of b] (c)
{\textbf{3. Pesaje y comparación}\\
Leer la cantidad acumulada\\
y comparar con el objetivo};

% Segunda fila
\node[etapa, below=0.65cm of c] (d)
{\textbf{4. Ajuste final}\\
Agregar harina remanente\\
sin un nuevo traslado};

\node[etapa, left=0.75cm of d] (e)
{\textbf{5. Devolución}\\
Retornar el recipiente\\
de transporte};

% Flujo principal
\draw[arr] (a) -- (b);
\draw[arr] (b) -- (c);
\draw[arr] (c) -- (d);
\draw[arr] (d) -- (e);

% Retorno para una nueva carga
\draw[arr]
  (c.north) -- ++(0,0.55) -|
  (b.north);
\node[anchor=south, align=center, text width=6cm, font=\footnotesize]
  at ($(c.north)+(-2.3,0.65)$)
  {Si se requiere otra carga principal};

\end{tikzpicture}

\par\vspace{0.2cm}

{\footnotesize
\textbf{Tara:} compensación de la masa del recipiente receptor
para iniciar el pesaje en cero.
}

\vfill

\raggedright
{\footnotesize Fuente: Elaboración propia con base en el procedimiento registrado.}

\end{frame}

% 6 — TRASLADOS
\begin{frame}{Traslados entre almacenamiento y punto de pesaje}
{\small\color{UCBBlue}\textbf{Etapa 2: traslado y cargas}}
\par\vspace{0.15cm}
\begin{columns}[c,onlytextwidth]
\begin{column}{0.72\textwidth}
\centering
\resizebox{\linewidth}{!}{% Planta esquemática: las coordenadas gráficas no representan una escala.
% El acceso y el giro de la puerta corresponden a la disposición indicada por el estudiante.
\begin{tikzpicture}[
    font=\small,
    equipo/.style={draw, line width=0.6pt, align=center, inner sep=5pt},
    relevante/.style={equipo, fill=black!5},
    cota/.style={thin, {Latex[length=4pt]}-{Latex[length=4pt]}},
    guia/.style={thin, densely dashed, black!60}
]
% Límite del área y abertura del acceso.
\draw[line width=0.8pt] (9,0) -- (0,0) -- (0,7) -- (14,7)
    -- (14,0) -- (10.3,0);
\draw[line width=0.6pt] (9,0) -- (9,1.3);
\draw[thin, dashed] (10.3,0) arc[start angle=0,end angle=90,radius=1.3];
\node[align=center, anchor=north] at (9.65,-0.18)
    {Acceso al área de trabajo};

% Equipos y superficies, en la disposición relativa del croquis original.
\node[equipo, minimum width=2cm, minimum height=2.1cm,
    text width=1.8cm, font=\footnotesize] at (1.3,5.65) {Almacén de\\ingredientes};
\node[equipo, minimum width=4.45cm, minimum height=1.5cm,
    text width=4cm] at (5.05,5.95) {Mesa de laminado\\y formado};
\node[relevante, minimum width=2.85cm, minimum height=1.6cm,
    text width=2.5cm] (harina) at (9.35,5.95) {Almacenamiento\\de harina};
\node[relevante, minimum width=2.25cm, minimum height=1.5cm,
    text width=1.85cm] at (12.55,5.95) {Amasadora};
\node[equipo, minimum width=4.8cm, minimum height=2.25cm,
    text width=4.25cm] at (6.15,3) {Mesa de reposo\\y porcionado};
\node[equipo, minimum width=1.6cm, minimum height=1.1cm,
    text width=1.25cm, font=\footnotesize, anchor=south] at (8,0.2) {Divisora\\de masa};

% La balanza se identifica dentro de la mesa vinculada a la dosificación.
\draw[fill=black!5, line width=0.6pt] (11.4,0.2) rectangle (13.7,4.6 );
\node[align=center, text width=2.05cm, font=\footnotesize] at (12.55,3.0)
    {Mesa de\\ingredientes y\\dosificación};
\node[equipo, fill=white, text width=1.65cm,
    minimum height=0.95cm, font=\footnotesize] (balanza) at (12.55,1.1)
    {Balanza\\Punto de pesaje};

% Cotas transcritas
\draw[guia] (10.2,5.15) -- (10.65,5.15);
\draw[guia] (10.2,1.1) -- (11.4,1.1);
\draw[cota] (10.3,1.1) -- node[left, fill=white, inner sep=2pt]
    {3,20 m} (10.3,5.15);
\draw[cota] (10.3,0.75) -- node[below, fill=white, inner sep=1pt,
    font=\footnotesize] {0,90 m} (11.4,0.75);
\draw[guia] (14,0) -- (14.7,0);
\draw[guia] (14,7) -- (14.7,7);
\draw[cota] (14.55,0) -- (14.55,7);
\node[rotate=90, anchor=north, inner sep=0pt] at (14.75,3.5) {4,10 m};

\end{tikzpicture}
}
\end{column}
\begin{column}{0.25\textwidth}
\centering
\includegraphics[width=\linewidth,height=0.34\textheight,keepaspectratio]
 {\marcoimg/recipiente_transporte.jpg}
\par\smallskip
{\small Recipiente de transporte}
\par\medskip
{\small El operario recoge, transporta y descarga harina en el recipiente receptor.}
\end{column}
\end{columns}
\vfill
{\footnotesize Fuente: Elaboración propia y registro fotográfico propio. Croquis sin escala.}
% Guion: señalar almacenamiento y balanza. No sumar cotas ni equiparar
% una cota del croquis a la longitud total recorrida. La distancia por trayecto
% requiere contraste entre PROJECT_CONTEXT.md y la sección 6 del registro.
\end{frame}

% 7 — PESAJE Y AJUSTE
\begin{frame}{Pesaje y ajuste manual de harina}
{\small\color{UCBBlue}\textbf{Etapas 3--5: pesaje, ajuste y comprobación}}
\par\vspace{0.15cm}
\begin{columns}[T,onlytextwidth]
\begin{column}{0.48\textwidth}
\centering
\includegraphics[width=\linewidth,height=0.43\textheight,keepaspectratio]
 {\marcoimg/balanza_recipiente_vacio.jpg}
\par\smallskip
{\small\bfseries Recipiente receptor sobre la balanza}
\par\smallskip
{\small La lectura se compara con la cantidad objetivo.}
\end{column}
\begin{column}{0.48\textwidth}
\centering
\includegraphics[width=\linewidth,height=0.43\textheight,keepaspectratio]
 {\marcoimg/ajuste_final_operario.jpg}
\par\smallskip
{\small\bfseries Ajuste final del operario}
\par\smallskip
{\small Se incorpora harina remanente y se comprueba la lectura final.}
\end{column}
\end{columns}
\vfill
{\footnotesize Fuente: Registro fotográfico propio y descripción del procedimiento observado.}
% Guion: distinguir el recipiente receptor del recipiente de transporte.
% La fotografía de la balanza ilustra el montaje, no una lectura final medida.
\end{frame}

% 8 — LEVANTAMIENTO DE CAMPO
\begin{frame}{Levantamiento de campo: siete ciclos registrados}
\begin{columns}[T,onlytextwidth]
\begin{column}{0.44\textwidth}
{\large\color{UCBBlue}\bfseries Jornadas documentadas}
\par\vspace{0.25cm}
{\Large\bfseries 02/09/2026}\par
Jornada ordinaria: \textbf{5 ciclos}
\par\vspace{0.4cm}
{\Large\bfseries 05/09/2026}\par
Sábado: \textbf{2 ciclos}
\par\vspace{0.35cm}
{\small Los resultados describen la muestra registrada.}
\end{column}
\begin{column}{0.51\textwidth}
{\large\color{UCBBlue}\bfseries Registro por ciclo}
\par\medskip
\begin{itemize}
\item Cantidad objetivo y cantidad final de harina.
\item Tiempo del procedimiento y actividades.
\end{itemize}
\medskip
\textbf{Intervalo analizado}\par\smallskip
{\small Desde la preparación del recipiente receptor hasta la devolución del recipiente de transporte.}
\end{column}
\end{columns}
\vfill
{\footnotesize Fuente: Registros de campo del 02/09 y 05/09/2026.}
% Guion: delimitar la unidad de análisis antes de mostrar los resultados.
% El ciclo sabatino de masa común cuenta una sola vez. A-003 sigue abierto:
% falta completar la segunda jornada ordinaria o documentar el cambio de plan.
% Pendiente: contraste con hojas originales e identificación del instrumento
% y procedimiento de cronometraje. No se atribuye un método no documentado.
\end{frame}

% BLOQUE 3 — EVIDENCIA Y PROBLEMA (9--12)
% 9 — RESULTADOS Y EVIDENCIA ASOCIADA
\begin{frame}{Resultados de los siete ciclos registrados}
\begin{columns}[T,onlytextwidth]
\begin{column}{0.47\textwidth}
{\color{UCBBlue}\Huge\bfseries 411,9 s}\par\smallskip
\textbf{Tiempo acumulado del procedimiento}
\par\vspace{0.35cm}
{\Large\bfseries 259,7 s\quad(63,0\%)}\par\smallskip
{\small Cargas principales y retornos por otra carga. Incluye recogida, transporte y descarga.}
\par\medskip
{\small Tres o cuatro cargas principales por ciclo.}
\end{column}
\begin{column}{0.47\textwidth}
{\color{UCBBlue}\Huge\bfseries +143 g}\par\smallskip
\textbf{Diferencia acumulada de harina}
\par\vspace{0.25cm}
\begin{tabular}{@{}lr@{}}
Cantidad objetivo & 48,000 kg\\[0.12cm]
Cantidad final & 48,143 kg
\end{tabular}
\par\medskip
{\Large\bfseries +11 a +38 g}\par\smallskip
{\small Diferencia final por ciclo. Los siete valores registrados fueron positivos.}
\end{column}
\end{columns}
\vfill
{\small\textbf{Cálculos derivados de los registros de campo.} La diferencia corresponde a cantidad final menos cantidad objetivo.}
\par\smallskip
{\footnotesize Fuente: Registros del 02/09 y 05/09/2026 y consolidación del diagnóstico.}
% Guion: describir tiempo e intervención, y luego diferencia de harina.
% El 63,0 % no representa caminata aislada. +143 g no equivale a desperdicio.
% Se conserva el rango observado; no constituye tolerancia de aceptación.
\end{frame}

% 10 — PROYECCIONES
\begin{frame}{Proyecciones a partir de las jornadas observadas}
{\small\textbf{Base registrada:} ordinaria, 304,5 s y +0,078 kg; sábado, 107,4 s y +0,065 kg.}
\par\vspace{0.45cm}
\begin{columns}[T,onlytextwidth]
\begin{column}{0.47\textwidth}
{\color{UCBBlue}\Large\bfseries Escenario semanal}\par\medskip
5 jornadas ordinarias + 1 sábado
\par\vspace{0.3cm}
{\color{UCBBlue}\LARGE\bfseries $\approx$ 27 min 10 s}\par\medskip
{\Large\bfseries +0,455 kg de harina}
\end{column}
\begin{column}{0.47\textwidth}
{\color{UCBBlue}\Large\bfseries Escenario mensual}\par\medskip
22 jornadas ordinarias + 4 sábados
\par\vspace{0.3cm}
{\color{UCBBlue}\LARGE\bfseries $\approx$ 1 h 59 min}\par\medskip
{\Large\bfseries +1,976 kg de harina}
\end{column}
\end{columns}
\vfill
{\small\textbf{Supuesto:} repetición de los ciclos y valores registrados en cada jornada. Los escenarios corresponden a proyecciones, no a mediciones semanales o mensuales.}
\par\medskip
{\footnotesize Fuente: Elaboración propia a partir de las jornadas registradas.}
% Guion: explicar la base y el supuesto antes del escenario mensual.
% Cálculos: 5*304,5+107,4=1629,9 s; 22*304,5+4*107,4=7128,6 s.
% Diferencias: 5*0,078+0,065=0,455 kg; 22*0,078+4*0,065=1,976 kg.
% La cobertura definitiva del relevamiento permanece abierta (A-003).
\end{frame}

% 11 — FORMULACIÓN
\begin{frame}{Formulación del problema}
\centering
{\small\color{UCBBlue}\textbf{Escenario mensual proyectado: 22 jornadas ordinarias y 4 sábados}}
\par\vspace{0.6cm}
\begin{minipage}{0.88\textwidth}
\raggedright\Large
En la Panadería San Miguel, el operario destina
\textbf{\color{UCBBlue}\mbox{1 h 59 min}}
al procedimiento manual de dosificación de harina, en el cual se acumula
una diferencia positiva de \textbf{\color{UCBBlue}1,976 kg}
respecto de las cantidades objetivo.
\end{minipage}
\vfill
\raggedright
{\small Valores derivados de las jornadas observadas.}
% Guion: leer como una única formulación. No afirmar ahorro, pérdida económica
% ni perjuicio de calidad. Se conserva literalmente la formulación de perfil.tex.
% D-003/D-004 y A-002: esta presentación no registra aprobación ni sustituye
% la versión académica de referencia por una nueva decisión.
\end{frame}

% 12 — ÁRBOL
% Se reemplaza el diagrama simplificado (una sola caja de "problema central")
% por la figura real del capítulo (images/marcoref/arbol_problemas.tex),
% reescalada. Motivo: la simplificación anterior fusionaba en un solo bloque
% los dos componentes que D-004 exige tratar conjuntamente (tiempo/intervención
% y diferencia de harina); la figura original los conserva como dos cajas
% enlazadas, preserva las subcausas y las bandas Efectos/Problema central/Causas,
% y evita mantener dos versiones del árbol que puedan desalinearse.
\begin{frame}{Árbol de problemas}
\centering
\resizebox{!}{0.68\textheight}{%
    \begin{minipage}{0.98\textwidth}
\centering

%%%%%%%%%%%%%%%%%%%%%%%%%%%%%%%%%%%%%%%%%%%%%%%%%%%%%%%%%%%%%%%%%%%%%%%%%%%%%%%
% EFECTOS
%%%%%%%%%%%%%%%%%%%%%%%%%%%%%%%%%%%%%%%%%%%%%%%%%%%%%%%%%%%%%%%%%%%%%%%%%%%%%%%

\textbf{EFECTOS OBSERVADOS Y PROYECTADOS}\\[0.5em]

\begin{tabular}{p{0.44\textwidth} p{0.44\textwidth}}

\fbox{\parbox{0.42\textwidth}{\centering
\textbf{E1}\\[0.3em]
Acumulación del tiempo destinado a la dosificación durante las jornadas
observadas\\[0.3em]
{\small
304,5 s en la jornada ordinaria\\
107,4 s en la jornada de sábado}
}}
&
\fbox{\parbox{0.42\textwidth}{\centering
\textbf{E2}\\[0.3em]
Acumulación mensual proyectada del tiempo destinado a la dosificación\\[0.3em]
{\small
7\,128,6 s para septiembre de 2026\\
1 h 58 min 48,6 s}
}}

\end{tabular}

\vspace{0.7em}

{\Large $\uparrow\qquad\qquad\qquad\uparrow$}

\vspace{0.7em}

%%%%%%%%%%%%%%%%%%%%%%%%%%%%%%%%%%%%%%%%%%%%%%%%%%%%%%%%%%%%%%%%%%%%%%%%%%%%%%%
% PROBLEMA CENTRAL
%%%%%%%%%%%%%%%%%%%%%%%%%%%%%%%%%%%%%%%%%%%%%%%%%%%%%%%%%%%%%%%%%%%%%%%%%%%%%%%

\fbox{\parbox{0.90\textwidth}{\centering
\textbf{PROBLEMA CENTRAL}\\[0.4em]

El procedimiento manual de dosificación de harina en la Panadería
San Miguel requiere intervención directa del operario durante
$53{,}2$--$66{,}6\ \mathrm{s}$ por ciclo.
}}

\vspace{0.7em}

{\Large $\uparrow\qquad\uparrow\qquad\uparrow$}

\vspace{0.7em}

%%%%%%%%%%%%%%%%%%%%%%%%%%%%%%%%%%%%%%%%%%%%%%%%%%%%%%%%%%%%%%%%%%%%%%%%%%%%%%%
% CAUSAS
%%%%%%%%%%%%%%%%%%%%%%%%%%%%%%%%%%%%%%%%%%%%%%%%%%%%%%%%%%%%%%%%%%%%%%%%%%%%%%%

\textbf{CAUSAS OBSERVADAS}\\[0.5em]

\begin{tabular}{p{0.30\textwidth} p{0.30\textwidth} p{0.30\textwidth}}

\fbox{\parbox{0.28\textwidth}{\centering
\textbf{C1}\\[0.3em]
Preparación del recipiente, aplicación de tara y comprobación de cero
en cada ciclo
}}
&
\fbox{\parbox{0.28\textwidth}{\centering
\textbf{C2}\\[0.3em]
Dosificación mediante cargas principales realizadas manualmente\\[0.3em]
{\small
3--4 cargas principales por ciclo}
}}
&
\fbox{\parbox{0.28\textwidth}{\centering
\textbf{C3}\\[0.3em]
Ajuste final y comprobación manual de la lectura de la balanza
}}

\end{tabular}

\end{minipage}}
\vfill\raggedright
{\footnotesize Fuente: Elaboración propia.}
% Guion: leer por bandas, de abajo hacia arriba (causas → problema central → efectos).
% Señalar el vínculo bidireccional entre los dos componentes del problema central
% (tiempo/intervención y diferencia de harina) como la revisión conjunta pedida
% por el docente (PROFESSOR_NOTES.md). No se interpreta la secuencia operativa
% como causalidad experimental demostrada. Figura idéntica a la de marcoref.tex:
% cualquier ajuste al árbol debe hacerse en el archivo fuente compartido,
% no de forma independiente en este slide.
\end{frame}

% BLOQUE 4 — RESPUESTA DEL PROYECTO (13--16)
\section{Objetivos}
% 13 — OBJETIVOS
\begin{frame}{Objetivo general y objetivos específicos}
{\large Desarrollar un \textbf{\color{UCBBlue}prototipo automatizado de dosificación gravimétrica de harina} para la Panadería San Miguel.}
\par\vspace{0.4cm}
\begin{columns}[T,onlytextwidth]
\begin{column}{0.48\textwidth}
{\small\color{UCBBlue}\textbf{Desarrollo}}
\begin{enumerate}
\item \textbf{Determinar} requerimientos funcionales y técnicos a partir del procedimiento manual.
\item \textbf{Diseñar} los subsistemas mecánico, de medición gravimétrica y de control.
\item \textbf{Implementar} el prototipo funcional conforme al diseño.
\end{enumerate}
\end{column}
\begin{column}{0.48\textwidth}
{\small\color{UCBBlue}\textbf{Verificación y evaluación}}
\begin{enumerate}\setcounter{enumi}{3}
\item \textbf{Verificar} los requerimientos mediante pruebas de funcionamiento.
\item \textbf{Comparar} el desempeño con el procedimiento manual bajo condiciones comparables.
\item \textbf{Determinar} el costo de construcción e integración mediante un desglose documentado.
\end{enumerate}
\end{column}
\end{columns}
\vfill
{\footnotesize Síntesis de los objetivos del documento de trabajo.}
% Guion: mantener implementación antes de verificación y comparación.
% Síntesis expositiva de marcoref.tex, no nueva formulación aprobada (A-002).
\end{frame}

% Motivación: no se añade contenido nuevo. La sección institucional permanece
% pendiente de desarrollo en el documento de trabajo.
\section{Justificación}
% 14 — JUSTIFICACIÓN TECNOLÓGICA
\begin{frame}{Justificación tecnológica}
\begin{columns}[T,onlytextwidth]
\begin{column}{0.47\textwidth}
{\color{UCBBlue}\large\bfseries Integración del prototipo}
\par\medskip
Diseño mecánico, medición gravimétrica, electrónica y control.
\par\vspace{0.35cm}
\textbf{Función prevista}\par\smallskip
Control de la alimentación a partir de la cantidad objetivo y del pesaje de la harina dosificada.
\end{column}
\begin{column}{0.47\textwidth}
{\color{UCBBlue}\large\bfseries Evaluación experimental}
\par\medskip
\begin{itemize}
\item Tiempo de dosificación.
\item Intervención del operario.
\item Diferencia respecto de la cantidad objetivo.
\item Repetibilidad.
\end{itemize}
\end{column}
\end{columns}
\vfill
{\small\textbf{Criterio de evaluación:} comparación con el procedimiento manual bajo condiciones comparables. Las mejoras requieren comprobación experimental.}
% Guion: vincular la función del sistema con las variables del diagnóstico.
% Los componentes, la estrategia de control y los umbrales siguen abiertos.
\end{frame}

% 15 — JUSTIFICACIÓN ECONÓMICA
\begin{frame}{Justificación económica: referencias comerciales}
\small
\renewcommand{\arraystretch}{1.35}
\begin{tabular}{@{}>{\raggedright\arraybackslash}p{0.22\textwidth}>{\raggedright\arraybackslash}p{0.34\textwidth}>{\raggedright\arraybackslash}p{0.36\textwidth}@{}}
\toprule
\textbf{Criterio} & \textbf{RUICHUAN} & \textbf{Trust-Long}\\
\midrule
Modelo & Serie SZC-T & TLP-500 / TLP-5000\\
Principio descrito & Dosificación gravimétrica por pérdida de peso & Operación por lotes y alimentación por tornillo\\
Precio de referencia & \textbf{USD 2\,000} & \textbf{USD 2\,600--3\,800}\\
\bottomrule
\end{tabular}
\par\vspace{0.4cm}
\textbf{\color{UCBBlue}Criterio para evaluar el bajo costo}\par\smallskip
Costo documentado de construcción e integración del prototipo frente a soluciones de funcionalidad comparable, considerando prestaciones y precio.
\vfill
{\footnotesize Precios conservados del documento de trabajo. Actualización y comparabilidad pendientes de verificación.}
\par\smallskip
{\footnotesize Fuente: \textcite{UnidadDosificadoraGravimetrica} y \textcite{MaquinaLlenadoPolvo}.}
% Guion: describir estas referencias como base de comparación, no demostrar
% bajo costo con los precios aislados. El costo del prototipo no está establecido.
% Las páginas comerciales no fueron accesibles durante esta revisión.
% A-001 sigue abierto. No se incorporan prestaciones ni cotizaciones nuevas.
\end{frame}

\section{Delimitación}
% 16 — ALCANCES Y LÍMITES
\begin{frame}{Alcances y límites del proyecto}
\begin{columns}[T,onlytextwidth]
\begin{column}{0.47\textwidth}
{\color{UCBBlue}\Large\bfseries Alcances}
\par\medskip
\small
\begin{itemize}
\item Dosificación gravimétrica automatizada de harina con cantidad objetivo configurable por el operario.
\item Construcción e integración del prototipo y pruebas de funcionamiento.
\item Comparación experimental con el procedimiento manual en San Miguel, bajo condiciones comparables.
\item Determinación del costo de construcción e integración.
\end{itemize}
\end{column}
\begin{column}{0.47\textwidth}
{\color{UCBBlue}\Large\bfseries Límites}
\par\medskip
\small
\begin{itemize}
\item Dosificación de otros ingredientes.
\item Mezclado y amasado, división y formado, fermentación y cocción.
\item Validación en otros establecimientos.
\item Análisis de rentabilidad o recuperación de la inversión.
\end{itemize}
\end{column}
\end{columns}
\vfill
{\footnotesize Rangos, capacidad, tolerancia, componentes y estrategia de control permanecen abiertos.}
% Guion: cerrar con función, caso de validación y exclusiones.
% Se conserva el alcance de perfil.tex. Los valores preliminares de marcoref.tex
% no se adoptan como requisitos, conforme a D-009 y A-004.
\end{frame}

% =========================================================
% TRANSICIÓN
% =========================================================

\begin{frame}[plain]

\vfill

\begin{center}

    {\color{UCBBlue}
    \fontsize{27}{31}\selectfont
    \bfseries
    Marco Referencial
    }

    \vspace{0.55cm}

    {\Large

    Contexto
    $\longrightarrow$
    Procedimiento
    $\longrightarrow$
    Evidencia

    \vspace{0.25cm}

    $\longrightarrow$

    \vspace{0.20cm}

    Problema
    $\longrightarrow$
    Objetivos
    $\longrightarrow$
    Delimitación

    }

\end{center}

\vfill

\end{frame}

% =========================================================
% PROPUESTA
% =========================================================

% La slide "Concepto de solución" se trasladaría a la
% sección posterior de Propuesta del perfil.
%
% No se incluye aquí para no mezclar la función del
% Marco Referencial con el Capítulo 3 del template.

\end{document}
